\chapter{Asymptotics}\label{ch:asymptotics}

\begin{center}\itshape
Divergent series are the invention of the devil, and it is shameful to base on them any demonstration whatsoever.\\
\upshape --- Niels Henrik Abel, letter to Holmboe (1826)
\end{center}

A million fractions are a long sum to add by hand. Four terms still write $H_{1000000}$ to the digits below:
\[
\ln10^6+\gamma+\frac1{2\cdot10^6}-\frac1{12\cdot10^{12}}=14.39272672286572363138\ldots
\]
Every digit on that line is a digit of $H_{1000000}$. The same four terms at $n=10$ give $2.9289674$, while $H_{10}=2.9289683$. Four terms, and the first six digits already agree. Where do $\frac1{2n}$ and $\frac1{12n^2}$ come from, what comes after them, and how many digits can you trust? That is the job of this chapter.

Many quantities in this volume have no closed form: $H_n$, $\ln n!$, the number of partitions of $n$, the cost of quicksort. What you can usually find is their \emph{size}: a simple function they approach, together with a bound on the error. Four tools do that work. The \emph{Euler--Maclaurin formula} compares a sum with an integral. The corrections are the Bernoulli numbers of Chapter~\ref{ch:finite}. Applied to $\ln\Gamma$, it gives \emph{Stirling's series}. The series diverges for every $x$, and from $x=4$ onward it still gives $\ln\Gamma(x)$ to twelve digits. \emph{Laplace's method} reads the size of $\int e^{n\varphi}$ from a neighbourhood of the maximum of $\varphi$. The last tool reads the coefficients of a generating function from its singularity, which is the promise made in Chapter~\ref{ch:generating}.

\portrait{abel}{Niels Henrik Abel (1802--1829). Portrait by Johan G\o rbitz, 1826. Public domain.}

Abel's sentence in the epigraph is from his letter to Holmboe of 16 January 1826. Divergent series had produced nonsense in careless hands, and Cauchy and Abel were putting analysis on definitions one could check. The rule in this chapter is the one they asked for. Every divergent series below comes with an explicit bound on its error, and the bound is what we use.

\begin{think}
Euler wrote $0!-1!+2!-3!+4!-\cdots\approx0.596$. The terms grow, so the series diverges. Can you find a function $f(x)$ whose Taylor coefficients at $0$ are $(-1)^kk!$, so that Euler's series is ``$f(1)$''? Try an integral: $k!=\int_0^\infty t^ke^{-t}\dd t$.
\end{think}

\section{Asymptotic expansions}

Let $f$ and $g$ be defined near $x_0$, a real number or $\pm\infty$, with $g$ non-zero near $x_0$. As $x\to x_0$ we write:
\begin{itemize}
\item $f=O(g)$ when $|f/g|$ is bounded near $x_0$;
\item $f=o(g)$ when $f/g\to0$;
\item $f\sim g$ when $f/g\to1$, that is, when $f=g+o(g)$.
\end{itemize}
These are the notations of Volume~1. An \emph{asymptotic expansion} asks for more than the first term.

\begin{definition}
Let $(\varphi_k)_{k\ge0}$ be functions defined near $x_0$ with $\varphi_{k+1}=o(\varphi_k)$, for instance $\varphi_k(x)=x^{-k}$ as $x\to\infty$ or $\varphi_k(x)=x^k$ as
$x\to0$. We say that $f$ has the \emph{asymptotic expansion} $\sum a_k\varphi_k$, and write $f\approx\sum_ka_k\varphi_k$, if for every $K\ge0$
\[
f-\sum_{k<K}a_k\varphi_k=O(\varphi_K).
\]
\end{definition}

The definition says nothing about convergence of $\sum a_k\varphi_k$. It fixes the number of terms $K$ and lets $x\to x_0$. Convergence fixes $x$ and lets $K\to\infty$. The two limits are different. The next statement shows that the first can hold when the second fails.

\begin{proposition}\label{prop:asymunique}
\begin{enumerate}[label=(\alph*)]
\item An asymptotic expansion is unique: if $f\approx\sum a_k\varphi_k$ and $f\approx\sum b_k\varphi_k$, then $a_k=b_k$ for every $k$.
\item Different functions can have the same expansion: $e^{-x}\approx\sum_k0\cdot x^{-k}$ as $x\to\infty$.
\end{enumerate}
\end{proposition}
\begin{proof}
\emph{Uniqueness.} Suppose $a_j=b_j$ for $j<K$. Subtract the two expansions at order $K+1$. Then $(a_K-b_K)\varphi_K=O(\varphi_{K+1})=o(\varphi_K)$, so $a_K-b_K=o(1)$. A constant that tends to $0$ is $0$, and $a_K=b_K$. Induction on $K$ finishes the argument.

\emph{The zero expansion.} For every $K$, $x^Ke^{-x}\to0$ (Volume~1, Chapter~8), so $e^{-x}=O(x^{-K})$.
\end{proof}

\begin{proposition}[Euler's series]\label{prop:eulerseries}
For $x>0$ let $f(x)=\displaystyle\int_0^\infty\frac{e^{-t}}{1+xt}\dd t$. For every $K\ge0$,
\[
\begin{gathered}f(x)=\sum_{k=0}^{K-1}(-1)^kk!\,x^k+r_K(x),\qq |r_K(x)|\le K!\,x^K,\end{gathered}
\]
and $r_K(x)$ has the sign $(-1)^K$ of the first omitted term. So $f(x)\approx\sum_k(-1)^kk!\,x^k$ as $x\to0^+$, although this series
diverges for every $x\ne0$.
\end{proposition}
\begin{proof}
\emph{The geometric identity.} For $u\ge0$ and $K\ge0$, summing a geometric progression gives $\frac1{1+u}=\sum_{k<K}(-u)^k+\frac{(-u)^K}{1+u}$.

\emph{Integrate.} Take $u=xt$, multiply by $e^{-t}$ and integrate. The formula $\int_0^\infty t^ke^{-t}\dd t=k!$ gives
\[
r_K(x)=(-1)^Kx^K\int_0^\infty\frac{t^Ke^{-t}}{1+xt}\dd t.
\]
The last integral is positive, and it is at most $\int_0^\infty t^Ke^{-t}\dd t=K!$, because $1+xt\ge1$.

\emph{Divergence.} The ratio of consecutive terms of the series is $(k+1)x$, which tends to infinity. The series diverges for every $x\ne0$.
\end{proof}

This answers the think box. Euler's ``$0!-1!+2!-\cdots$'' is $f(1)=\int_0^\infty\frac{e^{-t}}{1+t}\dd t=0.5963473623\ldots$

How do you use a series that diverges? Fix $x$. The bound $K!\,x^K$ decreases while $Kx<1$, and then it increases. The best choice is $K\approx\frac1x$. Stirling's formula turns that bound into $K!\,x^K\approx\sqrt{2\pi K}\,(Kx/e)^K\approx\sqrt{2\pi/x}\,e^{-1/x}$. At $x=0.1$, where $f(0.1)=0.9156333394\ldots$, the bound at $K=10$ is $3.6\cdot10^{-4}$ and the actual error is $1.77\cdot10^{-4}$. At $K=20$ the error has grown to $8\cdot10^{-3}$. This is \emph{optimal truncation}: stop at the smallest term. The error is then exponentially small in $\frac1x$. Any fixed number of terms does worse. The error never reaches $0$.

\begin{think}
In Proposition~\ref{prop:eulerseries} the error has the sign of the first omitted term, so $f(x)$ lies between two consecutive partial sums. Use the first eleven terms to trap $f(0.1)$ between two numbers you can certify.
\end{think}

\section{The Euler--Maclaurin formula}

Chapter~\ref{ch:finite} showed that $\sum_{k<n}k^p$ equals $\int_0^nx^p\dd x$ plus corrections built from Bernoulli numbers. The general statement replaces $x^p$ by a smooth function. The corrections use the \emph{periodic Bernoulli functions} $\bar B_m(x)=B_m(x-\lfloor x\rfloor)$, where the $B_m$ are the Bernoulli polynomials of Chapter~\ref{ch:finite}. Euler found the formula in the 1730s and used it to compute $\sum\frac1{n^2}$ to twenty decimal places, a few weeks before he found that the sum is $\frac{\pi^2}6$ (Sources, E47). Maclaurin published the formula independently in \emph{A Treatise of Fluxions} (Edinburgh, 1742).

\portrait{maclaurin}{Colin Maclaurin (1698--1746). Stipple engraving by S.~Freeman. Public domain.}

\begin{lemma}\label{lem:bernoullifn}
\begin{enumerate}[label=(\alph*)]
\item $B_m'=mB_{m-1}$ for $m\ge1$; $B_m(1)=B_m(0)=B_m$ for $m\ge2$; and $\int_0^1B_m(x)\dd x=0$ for $m\ge1$. So for $m\ge2$, $\bar B_m$ is continuous,
and $\bar B_m'=m\bar B_{m-1}$ off the integers.
\item For $m\ge2$ and every real $x$,
\[
\bar B_m(x)=-\frac{m!}{(2\pi i)^m}\sum_{n\ne0}\frac{e^{2\pi inx}}{n^m}.
\]
\item For $m\ge2$, $|\bar B_m(x)|\le\frac{2\,m!\,\zeta(m)}{(2\pi)^m}$. For $m=2k$ this bound is $|B_{2k}|$, and $(-1)^{k+1}\bigl(B_{2k}-\bar B_{2k}(x)\bigr)\ge0$,
with equality only at the integers.
\end{enumerate}
\end{lemma}
\begin{proof}
\emph{Derivatives and integrals.} Differentiate $B_m(x)=\sum_k\binom mkB_kx^{m-k}$ and use $(m-k)\binom mk=m\binom{m-1}k$. That is (a) for the derivative. Proposition~\ref{prop:bernoulli}(c) at $x=0$ gives $B_m(1)-B_m(0)=m\cdot0^{m-1}=0$ for $m\ge2$, and then $\int_0^1B_m=\frac{B_{m+1}(1)-B_{m+1}(0)}{m+1}=0$ for $m\ge1$. Continuity of $\bar B_m$ at an integer is the equality $B_m(1)=B_m(0)$.

\emph{Fourier coefficients.} Let $c_n(m)=\int_0^1B_m(x)e^{-2\pi inx}\dd x$, the Fourier coefficients of $\bar B_m$ for period $1$. By (a), $c_0(m)=0$ for $m\ge1$. Let $n\ne0$. For $m=1$, $B_1(x)=x-\frac12$ and an integration by parts gives $c_n(1)=-\frac1{2\pi in}$. For $m\ge2$, integrate by parts. The boundary term $\bigl[B_m(x)\frac{e^{-2\pi inx}}{-2\pi in}\bigr]_0^1$ vanishes because $B_m(1)=B_m(0)$, and $B_m'=mB_{m-1}$, so $c_n(m)=\frac m{2\pi in}c_n(m-1)$. Induction gives $c_n(m)=-\frac{m!}{(2\pi in)^m}$. For $m\ge2$, $\bar B_m$ is continuous and piecewise $C^1$, so its Fourier series converges to it at every point (Theorem~\ref{thm:dirichletfourier}, after $t=2\pi x$ makes the period $2\pi$).

\emph{The sign.} The bound in (c) follows from (b), since $\sum_{n\ne0}|n|^{-m}=2\zeta(m)$. For $m=2k$, $(2\pi i)^{2k}=(-1)^k(2\pi)^{2k}$. Group $n$ with $-n$:
\[
\bar B_{2k}(x)=(-1)^{k+1}\frac{2\,(2k)!}{(2\pi)^{2k}}\sum_{n\ge1}\frac{\cos2\pi nx}{n^{2k}}.
\]
At $x=0$ this is Euler's formula for $\zeta(2k)$ (Theorem~\ref{thm:zeta2k}), so the bound is $|B_{2k}|$. Subtracting,
$(-1)^{k+1}\bigl(B_{2k}-\bar B_{2k}(x)\bigr)=\frac{2\,(2k)!}{(2\pi)^{2k}}\sum_{n\ge1}\frac{1-\cos2\pi nx}{n^{2k}}\ge0$. The term $n=1$ is positive unless $x\in\Z$.
\end{proof}

\begin{theorem}[Euler--Maclaurin formula]\label{thm:eulermaclaurin}
Let $a<b$ be integers, $m\ge1$, and $f$ of class $C^m$ on $[a,b]$. Then
\[
\begin{aligned}\sum_{a<n\le b}f(n)&=\int_a^bf(x)\dd x+R_m\\&\quad+\sum_{j=1}^m\frac{(-1)^jB_j}{j!}\eqbr\quad\cdot\bigl(f^{(j-1)}(b)-f^{(j-1)}(a)\bigr),\end{aligned}
\]
with
\[
R_m=\frac{(-1)^{m+1}}{m!}\int_a^b\bar B_m(x)f^{(m)}(x)\dd x.
\]
\end{theorem}

Since $-B_1=\frac12$ and $B_j=0$ for odd $j\ge3$, the correction terms are $\frac{f(b)-f(a)}2$ and then $\frac{B_{2k}}{(2k)!}\bigl(f^{(2k-1)}(b)-f^{(2k-1)}(a)\bigr)$. With $B_2=\frac16$, $B_4=-\frac1{30}$ and $B_6=\frac1{42}$, the first three are $\frac1{12}(f'(b)-f'(a))$, $-\frac1{720}(f'''(b)-f'''(a))$ and $\frac1{30240}(f^{(5)}(b)-f^{(5)}(a))$.

\begin{proof}
\emph{The first order.} For $m=1$, on $]n,n+1[$ the periodic function is $\bar B_1(x)=x-n-\frac12$. Integration by parts gives
\[
\begin{aligned}\int_n^{n+1}\bar B_1f'&=\Bigl[\bigl(x-n-\tfrac12\bigr)f(x)\Bigr]_n^{n+1}-\int_n^{n+1}f\\&=\frac{f(n)+f(n+1)}2-\int_n^{n+1}f.\end{aligned}
\]
Sum over $a\le n<b$. The half-values add up to $\sum_{a<n\le b}f(n)-\frac{f(b)-f(a)}2$. That is the formula for $m=1$, because $(-1)^1B_1=\frac12$.

\emph{The inductive step.} Pass from $m$ to $m+1$, for $f$ of class $C^{m+1}$. The function $\bar B_{m+1}$ is continuous, its derivative is $(m+1)\bar B_m$ off the integers, and it equals $B_{m+1}$ at $a$ and at $b$ (Lemma~\ref{lem:bernoullifn}(a)). Integrate by parts on each $[n,n+1]$ and add:
\[
\begin{aligned}\int_a^b\bar B_mf^{(m)}&=\frac{B_{m+1}}{m+1}\bigl(f^{(m)}(b)-f^{(m)}(a)\bigr)\eqbr-\frac1{m+1}\int_a^b\bar B_{m+1}f^{(m+1)}.\end{aligned}
\]
Multiply by $\frac{(-1)^{m+1}}{m!}$:
\[
\begin{aligned}R_m&=\frac{(-1)^{m+1}B_{m+1}}{(m+1)!}\bigl(f^{(m)}(b)-f^{(m)}(a)\bigr)\eqbr+R_{m+1},\end{aligned}
\]
That is the formula at order $m+1$.
\end{proof}

The formula is an identity. The error has been moved into $R_m$, and $R_m$ is often small. Lemma~\ref{lem:bernoullifn}(c) bounds it.

\begin{corollary}[harmonic numbers]\label{cor:harmonic}
For every $M\ge1$, as $n\to\infty$,
\[
H_n=\ln n+\gamma+\frac1{2n}-\sum_{k=1}^{M-1}\frac{B_{2k}}{2k\,n^{2k}}+O\Bigl(\frac1{n^{2M}}\Bigr),
\]
and the error term has the sign of the first omitted term $-\frac{B_{2M}}{2M\,n^{2M}}$ and is smaller in absolute value. So $H_n=\ln n+\gamma+\frac1{2n}-\frac1{12n^2}+\frac1{120n^4}-\frac1{252n^6}+\cdots$
\end{corollary}
\begin{proof}
\emph{The formula at order $2M$.} Apply the theorem to $f(x)=\frac1x$ on $[1,n]$ at order $2M$. Here $f^{(2k-1)}(x)=-\frac{(2k-1)!}{x^{2k}}$ and $f^{(2M)}(x)=\frac{(2M)!}{x^{2M+1}}$, so
\[
\begin{aligned}H_n-1&=\ln n+\frac12\Bigl(\frac1n-1\Bigr)\eqbr+\sum_{k=1}^M\frac{B_{2k}}{2k}\Bigl(1-\frac1{n^{2k}}\Bigr)\eqbr-\int_1^n\frac{\bar B_{2M}(x)}{x^{2M+1}}\dd x.\end{aligned}
\]

\emph{The constant $\gamma$.} Since $\bar B_{2M}$ is bounded, the last integral converges as $n\to\infty$. Let $n\to\infty$ in $H_n-\ln n$, which tends to $\gamma$ (Volume~1, Chapter~3). Then $\gamma=\frac12+\sum_{k=1}^M\frac{B_{2k}}{2k}-\int_1^\infty\bar B_{2M}x^{-2M-1}\dd x$. Subtract:
\[
\begin{aligned}H_n&=\ln n+\gamma+\frac1{2n}-\sum_{k=1}^M\frac{B_{2k}}{2k\,n^{2k}}\eqbr+\int_n^\infty\frac{\bar B_{2M}(x)}{x^{2M+1}}\dd x.\end{aligned}
\]

\emph{The sign of the error.} Call $E_M(n)$ the error of the statement, the term $k=M$ of the sum together with the last integral, and call $U_M(n)=-\frac{B_{2M}}{2Mn^{2M}}$ the first omitted term. Since $\int_n^\infty x^{-2M-1}\dd x=\frac1{2Mn^{2M}}$,
\[
E_M(n)=\int_n^\infty\frac{\bar B_{2M}(x)-B_{2M}}{x^{2M+1}}\dd x.
\]
By Lemma~\ref{lem:bernoullifn}(c), $\bar B_{2M}-B_{2M}$ has the sign $(-1)^M$ and modulus at most $2|B_{2M}|$, and it is non-zero off the integers. So $E_M(n)$ has the sign $(-1)^M$ of $U_M(n)$, it is non-zero, and $|E_M(n)|\le\frac{|B_{2M}|}{Mn^{2M}}$. That is the $O$. The displayed formula holds for every $M$. Comparing it at $M$ and at $M+1$ gives $E_M=U_M+E_{M+1}$, and $E_{M+1}$ has the sign opposite to $U_M$. So $E_M-U_M$ has the sign opposite to $U_M$, while $E_M$ has the sign of $U_M$. Hence $0<E_M/U_M<1$.
\end{proof}

Figure~\ref{fig:rw11-harmonic} draws the error $H_n-\ln n-\gamma$ against the first three corrections. At $n=1$ the constant $\frac1{2n}$ is still $0.5$, visibly above the dot $0.423$. The term $-\frac1{12n^2}$ pulls the curve down to $0.417$, and $+\frac1{120n^4}$ puts it at $0.425$, next to the dot. By $n=8$ the three curves and the dots have met.

\begin{figure}[ht]
\centering
\fig{ch11-rewrite-harmonic}
\caption{The error $H_n-\ln n-\gamma$ (dots) against $\frac1{2n}$, then the same with $-\frac1{12n^2}$, then with $+\frac1{120n^4}$. Each correction pulls the curve onto the dots.}
\label{fig:rw11-harmonic}
\end{figure}

At $n=10$, the expansion through the term $-\frac1{252n^6}$ equals $2.928968253927\ldots$. The sum $H_{10}=\frac{7381}{2520}=2.928968253968\ldots$ is larger by $4.1\cdot10^{-11}$. The next term $\frac1{240n^8}$ is $4.2\cdot10^{-11}$. Ten terms of the sum and five terms of the expansion give eleven digits. The partial sum is the smaller number. The digits of $H_{10}$ are the larger one.

The same computation with $f(x)=x^{-s}$ continues $\zeta$ to the left, by a route of its own, alongside Chapter~\ref{ch:zeta}. For complex $s$ and $k\ge0$ write $(s)_k=s(s+1)\cdots(s+k-1)$, so that $\frac{\mathrm d^k}{\mathrm dx^k}x^{-s}=(-1)^k(s)_kx^{-s-k}$.

\begin{proposition}\label{prop:zetaem}
For integers $N\ge1$, $M\ge1$ and $\operatorname{Re}s>1$,
\[
\begin{aligned}\zeta(s)&=\sum_{n=1}^N\frac1{n^s}+\frac{N^{1-s}}{s-1}-\frac{N^{-s}}2\eqbr+\sum_{k=1}^M\frac{B_{2k}}{(2k)!}(s)_{2k-1}N^{1-s-2k}\eqbr-\frac{(s)_{2M}}{(2M)!}\int_N^\infty\frac{\bar B_{2M}(x)}{x^{s+2M}}\dd x.\end{aligned}
\]
The right side is holomorphic on $\operatorname{Re}s>1-2M$ except for a simple pole at $s=1$, so it is the continuation of $\zeta$ to that half-plane.
\end{proposition}
\begin{proof}
\emph{The tail.} Apply Theorem~\ref{thm:eulermaclaurin} to $f(x)=x^{-s}$ on $[N,b]$ at order $2M$ and let $b\to\infty$. For $\operatorname{Re}s>1$, every term at $b$ tends to $0$ and $\int_N^\infty x^{-s}\dd x=\frac{N^{1-s}}{s-1}$. The terms at $N$ are $-\frac{f(N)}2$ and $-\frac{B_{2k}}{(2k)!}f^{(2k-1)}(N)=\frac{B_{2k}}{(2k)!}(s)_{2k-1}N^{1-s-2k}$, and $R_{2M}=-\frac1{(2M)!}\int_N^\infty\bar B_{2M}(x)(s)_{2M}x^{-s-2M}\dd x$. Add $\sum_{n\le N}n^{-s}$.

\emph{Continuation.} On $\operatorname{Re}s\ge1-2M+\varepsilon$ the integrand is bounded by $|B_{2M}|x^{-1-\varepsilon}$, so the integral converges uniformly on compact sets of $\operatorname{Re}s>1-2M$ and is holomorphic there (Theorem~\ref{thm:holoparam}(b)). The other terms are entire, except $\frac{N^{1-s}}{s-1}$. The identity theorem (Theorem~\ref{thm:identity}) says this function agrees with the continuation of Chapter~\ref{ch:zeta} wherever both are defined.
\end{proof}

At $s=-n$ with $n\ge0$ and $2M\ge n+1$, the factor $(s)_{2M}=(-n)(1-n)\cdots(2M-1-n)$ contains the factor $0$, so the integral disappears and $\zeta(-n)$ is a finite combination of Bernoulli numbers. The two continuations agree, so the combination is the value found in Chapter~\ref{ch:zeta} (Theorem~\ref{thm:zetaneg} and Corollary~\ref{cor:zetaplane}). The smallest case is short enough to do by hand. Take $s=-1$ and $N=M=1$. The four terms are $1$, $\frac{N^2}{-2}=-\frac12$, $-\frac N2=-\frac12$ and $\frac{B_2}{2}\cdot(-1)=-\frac1{12}$. So $\zeta(-1)=1-\frac12-\frac12-\frac1{12}=-\frac1{12}$. No theta function entered the calculation. For $\operatorname{Re}s>1$ the same proposition is a practical way to compute $\zeta(s)$. Exercise~\ref{ex:zeta3} gets ten digits of $\zeta(3)$ from ten terms.

\section{Stirling's series}

Chapter~\ref{ch:gamma} proved $\Gamma(x)\sim\sqrt{2\pi}\,x^{x-\frac12}e^{-x}$ (Theorem~\ref{thm:stirlinggamma}). Euler--Maclaurin, applied to Gauss's product, gives the full expansion, and the sign of the remainder comes with it.

\begin{theorem}[Stirling's series]\label{thm:stirlingseries}
For $x>0$ and $m\ge1$,
\[
\begin{aligned}\ln\Gamma(x)&=\Bigl(x-\frac12\Bigr)\ln x-x+\frac12\ln2\pi\eqbr+\sum_{k=1}^{m-1}\frac{B_{2k}}{2k(2k-1)x^{2k-1}}+\rho_m(x),\end{aligned}
\]
where
\[
\rho_m(x)=\frac1{2m}\int_0^\infty\frac{B_{2m}-\bar B_{2m}(t)}{(x+t)^{2m}}\dd t.
\]
The remainder $\rho_m(x)$ has the sign of the first omitted term $T_m(x)=\frac{B_{2m}}{2m(2m-1)x^{2m-1}}$, and $|\rho_m(x)|<|T_m(x)|$.
\end{theorem}

With $B_2=\frac16$, $B_4=-\frac1{30}$, $B_6=\frac1{42}$, $B_8=-\frac1{30}$, the series reads
\[
\begin{aligned}\ln\Gamma(x)&\approx\Bigl(x-\frac12\Bigr)\ln x-x+\frac12\ln2\pi\eqbr+\frac1{12x}-\frac1{360x^3}+\frac1{1260x^5}-\frac1{1680x^7}+\cdots\end{aligned}
\]
For $m=1$ the theorem says $0<\rho_1(x)<\frac1{12x}$. For every integer $n\ge1$, apply it at $x=n$ and use $n!=n\Gamma(n)$:
$\sqrt{2\pi n}\,(n/e)^n<n!<\sqrt{2\pi n}\,(n/e)^ne^{1/(12n)}$.

\begin{proof}
\emph{Gauss's product.} By Gauss's formula (Chapter~\ref{ch:gamma}), $\ln\Gamma(x)=\lim_N\bigl(\ln N!+x\ln N-\sum_{k=0}^N\ln(x+k)\bigr)$. Apply Theorem~\ref{thm:eulermaclaurin} to $f(t)=\ln(x+t)$ on $[0,N]$ at order $2m$. Here $f^{(2k-1)}(t)=\frac{(2k-2)!}{(x+t)^{2k-1}}$ and $f^{(2m)}(t)=-\frac{(2m-1)!}{(x+t)^{2m}}$, and $\int_0^N\ln(x+t)\dd t=(x+N)\ln(x+N)-x\ln x-N$. Add the term $k=0$, which is $\ln x$:
\[
\begin{aligned}\sum_{k=0}^N\ln(x+k)&=\Bigl(x+N+\frac12\Bigr)\ln(x+N)-N\eqbr+\Bigl(\frac12-x\Bigr)\ln x\\
&\quad+\sum_{k=1}^m\frac{B_{2k}}{2k(2k-1)}\eqbr\cdot\Bigl(\frac1{(x+N)^{2k-1}}-\frac1{x^{2k-1}}\Bigr)\\
&\quad+\frac1{2m}\int_0^N\frac{\bar B_{2m}(t)}{(x+t)^{2m}}\dd t.\end{aligned}
\]
Let $N\to\infty$. The terms in $(x+N)^{1-2k}$ tend to $0$, and the integral converges absolutely. Also $(x+N+\frac12)\ln(x+N)=(x+N+\frac12)\ln N+x+o(1)$, because $\ln(x+N)-\ln N=\frac xN+O(N^{-2})$. Stirling's formula for $N!$ (Volume~1, Chapter~8) says $\ln N!-(N+\frac12)\ln N+N\to\frac12\ln2\pi$. Put these limits into Gauss's formula:
\[
\begin{aligned}\ln\Gamma(x)&=\Bigl(x-\frac12\Bigr)\ln x-x+\frac12\ln2\pi\eqbr+\sum_{k=1}^m\frac{B_{2k}}{2k(2k-1)x^{2k-1}}\eqbr-\frac1{2m}\int_0^\infty\frac{\bar B_{2m}(t)}{(x+t)^{2m}}\dd t.\end{aligned}
\]

\emph{Cancel the last term.} Since $\frac1{2m}\int_0^\infty\frac{B_{2m}}{(x+t)^{2m}}\dd t=\frac{B_{2m}}{2m(2m-1)x^{2m-1}}=T_m(x)$, the term $k=m$ of the sum and the integral combine into $\rho_m(x)$. That is the formula of the theorem.

\emph{The sign.} By Lemma~\ref{lem:bernoullifn}(c), $B_{2m}-\bar B_{2m}$ has the sign $(-1)^{m+1}$ of $B_{2m}$, and it is non-zero off the integers. So $\rho_m(x)$ has the sign of $T_m(x)$, and it is non-zero. The formula holds for every $m$. Comparing it at $m$ and at $m+1$ gives $\rho_m=T_m+\rho_{m+1}$. Now $\rho_{m+1}$ has the sign of $T_{m+1}$, which is opposite to the sign of $T_m$. So $\rho_m-T_m$ has the sign opposite to $T_m$, while $\rho_m$ has the sign of $T_m$. This forces $0<\rho_m/T_m<1$.
\end{proof}

\begin{figure}[ht]
\centering
\fig{ch11-stirling}
\caption{The error of Stirling's series for $\ln\Gamma(x)$, on a logarithmic scale, for $x=1,2,4$ and $m$ correction terms. The error falls, then grows without bound. The best $m$ is near $\pi x$, and the best errors are about $10^{-3.5}$, $10^{-6.4}$ and $10^{-12}$.}
\label{fig:stirlingerr}
\end{figure}

Since $|B_{2k}|=\frac{2(2k)!\,\zeta(2k)}{(2\pi)^{2k}}$, the terms $T_k(x)$ behave like $\frac{(2k-2)!}{(2\pi x)^{2k-1}}$, up to bounded factors. They decrease while $2k\lesssim2\pi x$, then grow faster than any geometric sequence. The series diverges for every $x$ (Figure~\ref{fig:stirlingerr}). Stirling's formula puts the smallest term, at $k\approx\pi x$, near $\frac{e^{-2\pi x}}{\pi\sqrt x}$, which is $2\cdot10^{-12}$ at $x=4$, as in the figure. We do not use this estimate, so the routine proof is left aside. The theorem already guarantees an error smaller than the first omitted term.

\begin{beauty}
The sign argument is the same one as in Proposition~\ref{prop:eulerseries}. The error lies between $0$ and the first omitted term, so consecutive partial sums of Stirling's series bracket $\ln\Gamma(x)$. You can certify digits without ever computing $\Gamma$. Tables of $\ln\Gamma$ were made this way for two centuries: Stirling's series at a large argument, then $\Gamma(x)=\Gamma(x+n)/\bigl(x(x+1)\cdots(x+n-1)\bigr)$ to come back down.
\end{beauty}

\section{Laplace's method}

Many quantities are integrals $\int e^{n\varphi(x)}\dd x$ with a large parameter $n$: $n!=\int_0^\infty t^ne^{-t}\dd t$, the Wallis integrals, the Beta function. For large $n$ the integrand is negligible away from the maximum of $\varphi$. Near that maximum, $\varphi$ looks like a downward parabola, so the integral looks like a Gaussian integral.

Figure~\ref{fig:rw11-bump} is the model case, the bumps $e^{-n x^2}$. Each has height $1$. The width shrinks like $1/\sqrt n$, and the area $\sqrt{\pi/n}$ shrinks the same way. That factor $1/\sqrt n$ is the one in Laplace's method. Figure~\ref{fig:laplace} is the same concentration for the integrand that produces $n!$.

\begin{figure}[ht]
\centering
\fig{ch11-rewrite-bump}
\caption{The bumps $e^{-nx^2}$ for $n=1$, $4$ and $16$. The height stays $1$. The width, and the area, shrink like $1/\sqrt n$.}
\label{fig:rw11-bump}
\end{figure}

\begin{figure}[ht]
\centering
\fig{ch11-laplace}
\caption{The integrands $(ue^{1-u})^n=e^{n(\ln u-u+1)}$ for $n=1,4,16,64$ (solid), and the Gaussians $e^{-n(u-1)^2/2}$ that match them near $u=1$ (dashed). The peak narrows like $1/\sqrt n$. Away from $u=1$ the integrands are exponentially small.}
\label{fig:laplace}
\end{figure}

\begin{theorem}[Laplace's method]\label{thm:laplace}
Let $I=\left]a,b\right[$ with $-\infty\le a<b\le+\infty$, let $\varphi:I\to\R$ be of class $C^2$ and $g:I\to\R$ continuous, with $\int_I|g|e^\varphi<\infty$.
Suppose that $c\in I$ satisfies $\varphi'(c)=0$, $\varphi''(c)<0$, and, for every $\delta>0$,
\[
\sup\bigl\{\varphi(x):x\in I,\ |x-c|\ge\delta\bigr\}<\varphi(c).
\]
Then, as $n\to\infty$,
\[
\begin{aligned}\int_Ig(x)e^{n\varphi(x)}\dd x&=e^{n\varphi(c)}\eqbr\cdot\biggl(g(c)\sqrt{\frac{2\pi}{n|\varphi''(c)|}}+o\Bigl(\frac1{\sqrt n}\Bigr)\biggr).\end{aligned}
\]
\end{theorem}

\begin{proof}
\emph{Normalise the maximum.} Replacing $\varphi$ by $\varphi-\varphi(c)$ multiplies both sides by $e^{-n\varphi(c)}$ and keeps the hypotheses, so assume $\varphi(c)=0$. Let $A=-\varphi''(c)>0$. We use $\int_\R e^{-\lambda u^2}\dd u=\sqrt{\pi/\lambda}$ for $\lambda>0$ (substitute $t=\lambda u^2$ and use $\Gamma(\frac12)=\sqrt\pi$, Proposition~\ref{prop:gammahalf}).

\emph{A parabola.} Fix $\varepsilon\in\left]0,\frac12\right[$, and choose $\delta>0$ with $[c-\delta,c+\delta]\subset I$ such that $|\varphi''(x)+A|\le\varepsilon A$ and $|g(x)-g(c)|\le\varepsilon$ for $|x-c|\le\delta$. Since $\varphi(c)=\varphi'(c)=0$, Taylor's formula with Lagrange remainder gives $\varphi(c+u)=\frac12\varphi''(\xi)u^2$ for some $\xi$ between $c$ and $c+u$, so for $|u|\le\delta$
\[
-(1+\varepsilon)\frac{Au^2}2\le\varphi(c+u)\le-(1-\varepsilon)\frac{Au^2}2.
\]

\emph{Near $c$.} Set $K_n=\int_{-\delta}^\delta e^{n\varphi(c+u)}\dd u$. The upper bound, with the Gaussian integral, gives $K_n\le\sqrt{2\pi/((1-\varepsilon)An)}$. The lower bound gives $K_n\ge\sqrt{2\pi/((1+\varepsilon)An)}-2\int_\delta^\infty e^{-\kappa nu^2}\dd u$ with $\kappa=\frac{(1+\varepsilon)A}2$. And $\int_\delta^\infty e^{-\kappa nu^2}\dd u\le\int_\delta^\infty\frac u\delta e^{-\kappa nu^2}\dd u=\frac{e^{-\kappa n\delta^2}}{2\kappa n\delta}$, which is $o(n^{-1/2})$. Also $\bigl|\int_{-\delta}^\delta g(c+u)e^{n\varphi(c+u)}\dd u-g(c)K_n\bigr|\le\varepsilon K_n$.

\emph{Away from $c$.} The hypothesis supplies $\eta>0$ with $\varphi(x)\le-\eta$ for $|x-c|\ge\delta$. There $e^{n\varphi}=e^{(n-1)\varphi}e^\varphi\le e^{-(n-1)\eta}e^\varphi$, so the integral of $ge^{n\varphi}$ over $\{|x-c|\ge\delta\}$ is at most $e^{-(n-1)\eta}\int_I|g|e^\varphi$, which is $o(n^{-1/2})$.

\emph{Conclusion.} Set $L=\sqrt{2\pi/A}$ and $J_n=\int_Ige^{n\varphi}$. The three estimates give
\[
\limsup_{n\to\infty}\bigl|\sqrt n\,J_n-g(c)L\bigr|\le L\,\eta(\varepsilon),
\]
where $\eta(\varepsilon)=|g(c)|\bigl((1-\varepsilon)^{-1/2}-1\bigr)+\varepsilon(1-\varepsilon)^{-1/2}$, and the comparison $1-(1+\varepsilon)^{-1/2}\le(1-\varepsilon)^{-1/2}-1$ was used. The right side tends to $0$ with $\varepsilon$. The left side does not depend on $\varepsilon$, so it is $0$.
\end{proof}

\begin{example}[Stirling's formula, a third proof]\label{ex:stirlinglaplace}
Substitute $t=nu$ in $n!=\int_0^\infty t^ne^{-t}\dd t$. Then $n!=n^{n+1}\int_0^\infty e^{n(\ln u-u)}\dd u$. Here $\varphi(u)=\ln u-u$ on $\left]0,\infty\right[$, $g=1$, and $\int_0^\infty e^\varphi=\int_0^\infty ue^{-u}\dd u=1$. The function $\varphi$ increases on $\left]0,1\right]$ and decreases on $[1,\infty[$, so the supremum off $]1-\delta,1+\delta[$ is $\max\bigl(\varphi(1-\delta),\varphi(1+\delta)\bigr)<\varphi(1)=-1$, and $\varphi''(1)=-1$. Laplace's method gives
\[
n!\sim n^{n+1}e^{-n}\sqrt{\frac{2\pi}n}=\sqrt{2\pi n}\Bigl(\frac ne\Bigr)^n.
\]
\end{example}

\begin{think}
In Figure~\ref{fig:laplace} the peak has width about $\frac1{\sqrt n}$. Where does the factor $\frac1{\sqrt n}$ in Laplace's method come from, and why does the formula use $\varphi''(c)$ and not $\varphi'''(c)$? What changes if the maximum sits at an endpoint $a$ with $\varphi'(a)<0$?
\end{think}

\section{The size of coefficients}

A generating function $F(z)=\sum a_nz^n$ that converges on a disc $|z|<R$ stores its coefficients in Cauchy's formula. For $0<r<R$ the series converges uniformly on the circle $|z|=r$. Integrate term by term and use $\int_{-\pi}^\pi e^{ik\theta}\dd\theta=2\pi[k=0]$:
\[
a_n=\frac1{2\pi}\int_{-\pi}^\pi\frac{F(re^{i\theta})}{r^ne^{in\theta}}\dd\theta.
\]
Two ideas read the size of $a_n$ from this integral. The first chooses $r$ carefully and applies Laplace's method. The second looks at the singularity of $F$ on the circle of convergence.

\subsection*{The saddle point}

\begin{example}[central binomial coefficients]\label{ex:centralsaddle}
Take $F(z)=(1+z)^{2n}$ and read off the coefficient of $z^n$. On the circle $|z|=r$, the integrand has modulus at most $\frac{(1+r)^{2n}}{r^n}$, and $\frac{(1+r)^2}r=r+2+\frac1r$ is smallest at $r=1$. On $|z|=1$, $(1+e^{i\theta})^{2n}e^{-in\theta}=\bigl(e^{i\theta/2}\cdot2\cos\frac\theta2\bigr)^{2n}e^{-in\theta}=\bigl(4\cos^2\frac\theta2\bigr)^n$, so
\[
\begin{gathered}\binom{2n}n=\frac1{2\pi}\int_{-\pi}^\pi e^{n\varphi(\theta)}\dd\theta,\qq\varphi(\theta)=\ln\Bigl(4\cos^2\frac\theta2\Bigr).\end{gathered}
\]
On $]-\pi,\pi[$, the function $\varphi$ is maximal only at $0$, with $\varphi(0)=\ln4$, and $\ln\cos^2\frac\theta2=-\frac{\theta^2}4+O(\theta^4)$, so $\varphi''(0)=-\frac12$. It increases on $]-\pi,0]$ and decreases on $[0,\pi[$, and $\int e^\varphi=\int4\cos^2\frac\theta2\dd\theta<\infty$. Laplace's method gives
\[
\binom{2n}n\sim\frac{4^n}{2\pi}\sqrt{\frac{2\pi}{n/2}}=\frac{4^n}{\sqrt{\pi n}}.
\]
\end{example}

The choice of $r$ is the whole idea. In general, choose $r$ to minimise $\frac{F(r)}{r^n}$. That is the point where $\frac{\mathrm d}{\mathrm dr}\bigl(\ln F(r)-n\ln r\bigr)=0$. As a function of the complex variable $z$, the modulus $|F(z)z^{-n}|$ then has a \emph{saddle point} at $z=r$: smallest along the real axis, largest along the circle. The contour crosses the pass at its lowest point, and Laplace's method does the rest. Chapter~\ref{ch:hr} takes the same idea to the partition function. The right radius there is $r=e^{-\pi/\sqrt{6n}}$, and the result is the size of $p(n)$.

\subsection*{Singularities}

When $F$ has radius of convergence $1$ and a singularity of the form $(1-z)^{-\alpha}$ at $z=1$, the coefficients grow like a power of $n$. The model case is an identity, not an estimate.

\begin{proposition}\label{prop:binomcoef}
Let $\alpha\in\R$, not a negative integer or $0$. The coefficients of $(1-z)^{-\alpha}=\sum_na_n(\alpha)z^n$ (for $|z|<1$) are
\[
a_n(\alpha)=\frac{\alpha(\alpha+1)\cdots(\alpha+n-1)}{n!}=\frac{\Gamma(n+\alpha)}{\Gamma(\alpha)\,n!},
\]
and, as $n\to\infty$,
\[
a_n(\alpha)=\frac{n^{\alpha-1}}{\Gamma(\alpha)}\Bigl(1+\frac{\alpha(\alpha-1)}{2n}+O\Bigl(\frac1{n^2}\Bigr)\Bigr).
\]
\end{proposition}
\begin{proof}
\emph{The series.} Let $S(z)=\sum_na_nz^n$ with $a_n$ the product formula. Since $\frac{a_{n+1}}{a_n}=\frac{\alpha+n}{n+1}\to1$, the series $S$ has radius of convergence $1$, and $(n+1)a_{n+1}-na_n=\alpha a_n$ says that $(1-z)S'(z)=\alpha S(z)$. With $(1-z)^\alpha=\exp\bigl(\alpha\operatorname{Log}(1-z)\bigr)$, holomorphic on the unit disc with derivative $-\alpha(1-z)^{\alpha-1}$, the derivative of $(1-z)^\alpha S(z)$ is $(1-z)^{\alpha-1}\bigl((1-z)S'-\alpha S\bigr)=0$. So $(1-z)^\alpha S(z)$ is constant, equal to its value $1$ at $0$. The Gamma form follows from $\Gamma(y+1)=y\Gamma(y)$, applied $n$ times.

\emph{Stirling at two points.} For the asymptotics, apply Theorem~\ref{thm:stirlingseries} with $m=2$ at $y=n+\alpha$ and at $y=n+1$ (both positive for $n$ large): $\ln\Gamma(y)=(y-\frac12)\ln y-y+\frac12\ln2\pi+\frac1{12y}+\rho_2(y)$ with $|\rho_2(y)|<\frac1{360y^3}$. The difference of the terms $\frac1{12y}+\rho_2(y)$ at the two points is $O(n^{-2})$. With $\ln(1+\frac\beta n)=\frac\beta n-\frac{\beta^2}{2n^2}+O(n^{-3})$,
\[
\begin{aligned}\Bigl(n+\beta-\frac12\Bigr)\ln\Bigl(1+\frac\beta n\Bigr)&=\beta+\frac{\beta(\beta-1)}{2n}\eqbr+O\Bigl(\frac1{n^2}\Bigr),\end{aligned}
\]
and $\ln(n+\beta)=\ln n+\ln(1+\frac\beta n)$. Subtract the case $\beta=1$ from the case $\beta=\alpha$, and subtract the terms $-(n+\alpha)+(n+1)$. Then $\ln\Gamma(n+\alpha)-\ln\Gamma(n+1)=(\alpha-1)\ln n+\frac{\alpha(\alpha-1)}{2n}+O(n^{-2})$. Exponentiate.
\end{proof}

For $\alpha=\frac12$, the series $\sum_n\binom{2n}nz^n$ equals $(1-4z)^{-1/2}$, and the proposition gives $\binom{2n}n=\frac{4^n}{\sqrt{\pi n}}\bigl(1-\frac1{8n}+O(n^{-2})\bigr)$. That is a second proof of Example~\ref{ex:centralsaddle}, with one more term. The Catalan numbers $\frac1{n+1}\binom{2n}n$ are therefore $\frac{4^n}{\sqrt\pi\,n^{3/2}}\bigl(1+O(\frac1n)\bigr)$.

The next theorem says that the model case controls every function with the same singularity, once you multiply by a factor that stays analytic across the circle.

\begin{theorem}[transfer]\label{thm:transfer}
Let $\alpha\in\R$, not a negative integer or $0$, and let $h(z)=\sum_kh_kz^k$ be a power series with radius of convergence greater than $1$. Then the
coefficients of $F(z)=(1-z)^{-\alpha}h(z)$ satisfy
\[
[z^n]F(z)=\frac{h(1)}{\Gamma(\alpha)}n^{\alpha-1}+O(n^{\alpha-2}).
\]
\end{theorem}

\begin{proof}
\emph{Split the convolution.} Fix $\rho>1$ less than the radius of $h$. Then $h_k\rho^k\to0$, so $|h_k|\le C\rho^{-k}$ for some $C$, and $\sum_kk|h_k|<\infty$. Write $a_j=a_j(\alpha)$, so that $[z^n]F=\sum_{k=0}^nh_ka_{n-k}$ (the product of the power series) and $h(1)a_n=\sum_{k\ge0}h_ka_n$. Then
\[
\begin{aligned}{}[z^n]F-h(1)a_n&=\sum_{k\le n/2}h_k(a_{n-k}-a_n)\eqbr+\sum_{n/2<k\le n}h_ka_{n-k}-\sum_{k>n/2}h_ka_n.\end{aligned}
\]
\emph{The near terms.} Since $(1-z)\cdot(1-z)^{-\alpha}=(1-z)^{-(\alpha-1)}$, the differences $a_j-a_{j-1}$ are the coefficients of $(1-z)^{-(\alpha-1)}$. If $\alpha-1$ is not a negative integer or $0$, they are $O(j^{\alpha-2})$ by Proposition~\ref{prop:binomcoef}. Otherwise $(1-z)^{1-\alpha}$ is a polynomial and the differences vanish for large $j$. In both cases, for $k\le\frac n2$, $|a_n-a_{n-k}|\le\sum_{n-k<j\le n}|a_j-a_{j-1}|\le C'kn^{\alpha-2}$, where $C'$ depends only on $\alpha$. The indices $j$ lie between $\frac n2$ and $n$. So the first sum is at most $C'n^{\alpha-2}\sum_kk|h_k|$.

\emph{The far terms.} Proposition~\ref{prop:binomcoef} also gives $|a_j|\le C''(1+j)^{|\alpha|+1}$ for all $j$, so the last two sums are at most $2C''(1+n)^{|\alpha|+1}\sum_{k>n/2}C\rho^{-k}$, which decreases exponentially. Hence $[z^n]F=h(1)a_n+O(n^{\alpha-2})$, and $a_n=\frac{n^{\alpha-1}}{\Gamma(\alpha)}+O(n^{\alpha-2})$ by Proposition~\ref{prop:binomcoef}.
\end{proof}

\begin{example}[2-regular graphs]\label{ex:tworegular}
A graph on $\{1,\dots,n\}$ in which every vertex has degree $2$ is a disjoint union of cycles of length at least $3$. On $k\ge3$ labelled vertices there
are $\frac{(k-1)!}2$ cycles (orderings up to rotation and reflection). By the exponential formula (Theorem~\ref{thm:expformula}), the number $g_n$ of such graphs
satisfies
\[
\sum_{n\ge0}g_n\frac{z^n}{n!}=\exp\Bigl(\sum_{k\ge3}\frac{z^k}{2k}\Bigr)=\frac{e^{-z/2-z^2/4}}{\sqrt{1-z}},
\]
since $\sum_{k\ge1}\frac{z^k}k=-\ln(1-z)$. Here $\alpha=\frac12$ and $h(z)=e^{-z/2-z^2/4}$ is entire, with $h(1)=e^{-3/4}$. So
\[
\frac{g_n}{n!}=\frac{e^{-3/4}}{\sqrt{\pi n}}+O\bigl(n^{-3/2}\bigr).
\]
The first values are $g_3=1$, $g_4=3$, $g_5=12$, $g_6=70$, $g_7=465$. At $n=30$ the approximation is $2\%$ too large, which is the size the $O(n^{-3/2})$ predicts.
\end{example}

The same reasoning gives the derangements. A permutation with no fixed point is a set of cycles of length at least $2$, so its exponential generating function is $\exp\bigl(-\ln(1-z)-z\bigr)=\frac{e^{-z}}{1-z}$. The theorem with $\alpha=1$ gives $\frac{D_n}{n!}=e^{-1}+O(\frac1n)$. The truth is sharper: $\frac{D_n}{n!}=\sum_{k\le n}\frac{(-1)^k}{k!}$ differs from $e^{-1}$ by less than $\frac1{(n+1)!}$. The error term in the theorem is only an upper bound. It is sharp when $h$ fails to continue far past the circle, and it is loose for an entire $h$ with $\alpha=1$, where every coefficient of $h(z)/(1-z)$ is a partial sum of $h(1)$.

\begin{remark}
Flajolet and Odlyzko (1990) proved the transfer theorem under weaker hypotheses. The function $F$ needs only to be analytic in a domain shaped like a disc with a bite taken at $z=1$, with $F(z)\sim C(1-z)^{-\alpha}$ near $1$, and logarithmic factors are allowed. Their proof uses a Cauchy integral on a contour that stays close to the singularity. It is the main tool of analytic combinatorics; see the book of Flajolet and Sedgewick.
\end{remark}

\section{What we proved, and what comes next}

The four terms of the opening are Corollary~\ref{cor:harmonic}. The Euler--Maclaurin formula (Theorem~\ref{thm:eulermaclaurin}) turns $H_n$ into $\ln n+\gamma+\frac1{2n}-\frac1{12n^2}+\cdots$, and the error is smaller than the first term left out. For $n=10^6$ that term is $\frac1{120n^4}\approx10^{-26}$. The same formula continued $\zeta$ to the left by a second route (Proposition~\ref{prop:zetaem}) and gave Stirling's series with the sign of the remainder (Theorem~\ref{thm:stirlingseries}). Laplace's method (Theorem~\ref{thm:laplace}) read the size of an integral from its peak, and the transfer theorem (Theorem~\ref{thm:transfer}) read the size of coefficients from a singularity. A closed expansion is a piece of order. The bound on the error, which Abel asked for, is what makes the digits safe to print.

Two chapters remain. Chapter~\ref{ch:hr} puts the saddle point together with the symmetry of Chapter~\ref{ch:zeta}, and asks how large $p(n)$ is. Exercise~\ref{ex:partitionbound} is the first step of that argument. Chapter~\ref{ch:agm} computes $\pi$ by a sequence that doubles its correct digits at every step. That is the opposite behaviour of an asymptotic series.

\section{Exercises}

\begin{exercise}[$\star$]\label{ex:kk}
Show that $\sum_{k=1}^nk^k=n^n\bigl(1+\frac1{en}+O(\frac1{n^2})\bigr)$.
\end{exercise}

\begin{exercise}[$\star$]\label{ex:wallislaplace}
Use Laplace's method to show that $W_n=\int_0^{\pi/2}\cos^n\theta\dd\theta\sim\sqrt{\frac\pi{2n}}$.
\end{exercise}

\begin{exercise}[$\star\star$]\label{ex:quickexp}
Example~\ref{ex:quicksortexact} of Chapter~\ref{ch:generating} showed that randomised quicksort makes on average $C_n=2(n+1)H_n-4n$ comparisons. Prove
\[
\begin{aligned}C_n&=2n\ln n+(2\gamma-4)n+2\ln n\eqbr+2\gamma+1+\frac5{6n}+O\Bigl(\frac1{n^2}\Bigr).\end{aligned}
\]
\end{exercise}

\begin{exercise}[$\star\star$]\label{ex:zeta3}
Apply Proposition~\ref{prop:zetaem} with $s=3$, $N=10$ and the terms up to $B_6$. Prove
\[
\begin{aligned}\zeta(3)&=\sum_{n=1}^{10}\frac1{n^3}+\frac1{200}-\frac1{2000}\eqbr+\frac1{4\cdot10^4}-\frac1{12\cdot10^6}+\frac1{12\cdot10^8}+r,\end{aligned}
\]
with $|r|\le1.1\cdot10^{-10}$. Use the remainder of order $7$ rather than $6$. Compute $\zeta(3)$ to ten digits.
\end{exercise}

\begin{exercise}[$\star\star$]\label{ex:partitionbound}
Let $p(n)$ be the number of partitions of $n$, and let $P(x)=\sum_np(n)x^n=\prod_{k\ge1}\frac1{1-x^k}$ for $0<x<1$ (Chapter~\ref{ch:generating}).
\begin{enumerate}[label=(\alph*)]
\item Prove $\ln P(e^{-t})\le\frac{\pi^2}{6t}$ for $t>0$.
\item Deduce $\ln p(n)\le\pi\sqrt{2n/3}$ for every $n\ge1$.
\item Prove $p(n)\ge2^{\lfloor\sqrt n\rfloor-1}$, so that $\ln p(n)$ is of order exactly $\sqrt n$.
\end{enumerate}
\end{exercise}

\begin{exercise}[$\star\star$]\label{ex:motzkinasym}
The Motzkin numbers of Exercise~\ref{ex:motzkin} have generating function $M(z)=\frac{1-z-\sqrt{(1-3z)(1+z)}}{2z^2}$. Prove
\[
M_n\sim\frac{3^{n+3/2}}{2\sqrt\pi\,n^{3/2}}.
\]
\end{exercise}

\begin{exercise}[$\star\star\star$]\label{pb:unlabelled}
Let $g_n$ be the number of graphs on $n$ vertices up to isomorphism, and set $N=\binom n2$.
\begin{enumerate}[label=(\alph*)]
\item Prove $g_n\ge\frac{2^N}{n!}$.
\item Prove $g_n\sim\frac{2^N}{n!}$. Almost all graphs have no symmetry.
\item Prove the sharper form $\frac{n!\,g_n}{2^N}=1+\frac{n(n-1)}{2^{n-1}}+O\bigl(n^42^{-3n/2}\bigr)$.
\end{enumerate}
\end{exercise}

\begin{exercise}[$\star\star$]\label{pb:oddcycles}
Let $o_n$ count the permutations of $\{1,\dots,n\}$ whose cycles all have odd length, and let $e_n$ count those whose cycles all have even length.
\begin{enumerate}[label=(\alph*)]
\item Prove $o_n=e_n=\bigl(1\cdot3\cdot5\cdots(n-1)\bigr)^2$ for $n$ even, and $o_n=n\,o_{n-1}$ for $n$ odd.
\item Deduce that the probability a random permutation has odd order is asymptotic to $\sqrt{\frac2{\pi n}}$.
\end{enumerate}
\end{exercise}

\begin{exercise}[$\star$]\label{pb:powersum}
\emph{The last terms win.} Prove
\[
\frac1{n^n}\sum_{k=1}^nk^n=\frac e{e-1}-\frac{e(e+1)}{2(e-1)^3}\cdot\frac1n+O\Bigl(\frac1{n^2}\Bigr).
\]
Write the sum as $\sum_{j=0}^{n-1}(1-\frac jn)^n$. For $n=100$ the left side is $1.57211$ and the two-term approximation is $1.57202$.
\end{exercise}

\begin{exercise}[$\star\star$]\label{pb:birthdaytriple}
\emph{Three people with the same birthday, for large $n$.} Draw uniformly and independently from $n$ values, and let $T$ be the number of draws until some value has appeared three times. Exercise~\ref{pb:poissonise} of Chapter~\ref{ch:generating} shows that
\[
\mathbb E[T]=\int_0^\infty e^{-t}\Bigl(1+\frac tn+\frac{t^2}{2n^2}\Bigr)^n\dd t.
\]
\begin{enumerate}[label=(\alph*)]
\item Prove $\mathbb E[T]\sim\Gamma(\frac43)(6n^2)^{1/3}$ as $n\to\infty$.
\item For $n=365$ the exact value is $88.74$. What does (a) give?
\end{enumerate}
\end{exercise}

\begin{exercise}[$\star\star$]\label{pb:sqrtsum}
\emph{The constant in $1+\sqrt2+\dots+\sqrt n$.} Use the continuation of $\zeta$ by Euler--Maclaurin (Proposition~\ref{prop:zetaem}) to prove
\[
\begin{aligned}\sum_{k=1}^n\sqrt k&=\frac23n^{3/2}+\frac12\sqrt n+\zeta\Bigl(-\frac12\Bigr)+\frac1{24\sqrt n}+O\bigl(n^{-5/2}\bigr),\\
\sum_{k=1}^n\frac1{\sqrt k}&=2\sqrt n+\zeta\Bigl(\frac12\Bigr)+\frac1{2\sqrt n}-\frac1{24n^{3/2}}+O\bigl(n^{-7/2}\bigr),\end{aligned}
\]
where $\zeta(-\frac12)=-0.2078862\ldots$ and $\zeta(\frac12)=-1.4603545\ldots$ So the constant left over when the growing terms are removed is a value of the zeta function.
\end{exercise}

\begin{exercise}[$\star$]\label{ex:rw11a}
In Proposition~\ref{prop:eulerseries} take $x=\frac15$. Show that the partial sum through $k=3$ equals $\frac{104}{125}$, and that $f\bigl(\frac15\bigr)$ lies strictly between $\frac{104}{125}$ and $\frac{544}{625}$.
\end{exercise}

\begin{exercise}[$\star\star$]\label{ex:rw11b}
Use Corollary~\ref{cor:harmonic} through the term $\frac1{120n^4}$ at $n=5$. Compare the approximation with $H_5=\frac{137}{60}$, and check that the difference is smaller than $\bigl|\frac1{252\cdot5^6}\bigr|$ and has the sign of $-\frac1{252n^6}$.
\end{exercise}

\begin{exercise}[$\star\star$]\label{ex:rw11c}
The maximum of $\varphi(x)=-x$ on $[0,\infty)$ sits at the endpoint $0$, where $\varphi'(0)=-1$, so Theorem~\ref{thm:laplace} does not apply. Show that $\int_0^\infty e^{-nx}\dd x=\frac1n$. For comparison, $\varphi(x)=-x^2$ on $\R$ does meet the hypotheses, and $\int_{-\infty}^\infty e^{-nx^2}\dd x=\sqrt{\pi/n}$. At $n=10$, compare the two values. Figure~\ref{fig:rw11-bump} is the second integrand.
\end{exercise}

\begin{exercise}[$\star\star\star$]\label{ex:rw11d}
Example~\ref{ex:centralsaddle} gives $\binom{2n}{n}\sim 4^n/\sqrt{\pi n}$, and Proposition~\ref{prop:binomcoef} adds one term:
\[
\binom{2n}{n}=\frac{4^n}{\sqrt{\pi n}}\Bigl(1-\frac1{8n}+O(n^{-2})\Bigr).
\]
For $n=2$, compare $\binom42=6$ with $4^2/\sqrt{2\pi}$. For $n=10$, compare $\binom{20}{10}=184756$ with the leading term and with the two-term approximation.
\end{exercise}

\section{Hints}

\begin{hint}{ex:kk}
The last term is $n^n$. The term before it is $n^{n-1}(1-\frac1n)^{n-1}$. Bound the earlier terms together, crudely.
\end{hint}
\begin{hint}{ex:wallislaplace}
Use the symmetry $\theta\mapsto-\theta$ so that the maximum of $\ln\cos\theta$ sits inside the interval.
\end{hint}
\begin{hint}{ex:quickexp}
Use Corollary~\ref{cor:harmonic} with two correction terms. Multiply out and collect powers of $n$.
\end{hint}
\begin{hint}{ex:zeta3}
Compute $(3)_{2k-1}$. For the remainder, $R_6=R_7$ because $B_7=0$, and Lemma~\ref{lem:bernoullifn}(c) bounds $\bar B_7$.
\end{hint}
\begin{hint}{ex:partitionbound}
(a) Expand $-\ln(1-y)=\sum_m\frac{y^m}m$ and swap the sums. Use $e^u-1\ge u$. (b) Start from $p(n)x^n\le P(x)$ and optimise in $t$. (c) Take parts from $\{2,\dots,\lfloor\sqrt n\rfloor\}$ and fill the rest with $1$'s.
\end{hint}
\begin{hint}{ex:motzkinasym}
Read the coefficient of $z^{n+2}$ in $\sqrt{(1-3z)(1+z)}$. Set $w=3z$ to move the singularity to $1$.
\end{hint}

\begin{hint}{pb:unlabelled}
Use Burnside, and set the identity term apart. A permutation that moves $m$ points leaves many pairs of vertices outside its fixed pairs. The pairs that move come in cycles of length at least $2$.
\end{hint}

\begin{hint}{pb:oddcycles}
The generating functions are $\sqrt{\frac{1+x}{1-x}}$ and $\frac1{\sqrt{1-x^2}}$.
\end{hint}

\begin{hint}{pb:powersum}
For $0\le x\le\frac12$, one has $-x-\frac{x^2}2-x^3\le\ln(1-x)\le-x-\frac{x^2}2$. Use that for $j\le\frac n2$, and bound the terms with $j>\frac n2$ by $2^{-n}$. You will need $\sum_jj^2x^j=\frac{x(1+x)}{(1-x)^3}$.
\end{hint}

\begin{hint}{pb:birthdaytriple}
Substitute $t=nx$. The function $\psi(x)=e^{-x}(1+x+\frac{x^2}2)$ decreases, and $\ln\psi(x)=-\frac{x^3}6+O(x^4)$ near $0$. The maximum is at the endpoint $0$, and it is flat to third order, so Theorem~\ref{thm:laplace} does not apply as stated. Its proof still does: squeeze $\psi$ between $e^{-(1\pm\varepsilon)x^3/6}$ near $0$.
\end{hint}

\begin{hint}{pb:sqrtsum}
In Proposition~\ref{prop:zetaem}, take $s=-\frac12$ (and, separately, $s=\frac12$) with $M=2$, and replace $N$ by $n$. Bound the integral remainder.
\end{hint}


\begin{hint}{ex:rw11a}
The remainder after three terms has the sign of $(-1)^4$ and is strictly smaller than $4!\,(1/5)^4$.
\end{hint}

\begin{hint}{ex:rw11b}
The first omitted term after $\frac1{120n^4}$ is $-\frac1{252n^6}$. Corollary~\ref{cor:harmonic} says the error has that sign and is smaller in absolute value.
\end{hint}

\begin{hint}{ex:rw11c}
The exponential integral is the antiderivative $-e^{-nx}/n$. For the Gaussian, set $u=x\sqrt n$ and use $\int_{-\infty}^\infty e^{-u^2}\dd u=\sqrt\pi$.
\end{hint}

\begin{hint}{ex:rw11d}
The leading term is $4^n/\sqrt{\pi n}$. Multiply by $1-1/(8n)$ for the second comparison. For $n=10$ the binomial coefficient is $184756$.
\end{hint}

\section{Solutions}

\begin{solution}{ex:kk}
\idea In $\sum k^k$ the last term is so large that the sum is essentially its last two terms. The only job is to bound the rest.
\writeup The last two terms are $n^n$ and $(n-1)^{n-1}=n^{n-1}(1-\frac1n)^{n-1}$. Now $(n-1)\ln(1-\frac1n)=-1+\frac1{2n}+O(n^{-2})$, so $(1-\frac1n)^{n-1}=e^{-1}+O(\frac1n)$, and $(n-1)^{n-1}=\frac{n^n}{en}+O(n^{n-2})$. The term $k=n-2$ is at most $n^{n-2}$. Each of the at most $n$ terms with $k\le n-3$ is at most $(n-3)^{n-3}\le n^{n-3}$, so together they are at most $n^{n-2}$. Altogether $\sum_{k\le n}k^k=n^n+\frac{n^n}{en}+O(n^{n-2})$. (At $n=100$ the $O$ term divided by $n^{n-2}$ is about $0.32$.)
\end{solution}

\begin{solution}{ex:wallislaplace}
\idea The integral $\cos^n\theta=e^{n\ln\cos\theta}$ is a Laplace integral. The maximum $\theta=0$ is an endpoint of $[0,\frac\pi2]$, and symmetry moves it inside.
\writeup Because $\cos$ is even,
\[
\begin{gathered}W_n=\frac12\int_{-\pi/2}^{\pi/2}e^{n\varphi(\theta)}\dd\theta,\qq\varphi=\ln\cos.\end{gathered}
\]
Then $\varphi'=-\tan$ and $\varphi''=-1-\tan^2$, so $\varphi(0)=0$, $\varphi'(0)=0$, $\varphi''(0)=-1$. The function $\varphi$ increases up to $0$ and decreases after, so its supremum for $|\theta|\ge\delta$ is $\ln\cos\delta<0$. And $\int e^\varphi=\int\cos=2$. Theorem~\ref{thm:laplace} with $g=1$ gives $W_n\sim\frac12\sqrt{2\pi/n}=\sqrt{\frac\pi{2n}}$. Combined with $W_{2n}W_{2n+1}$, computed by integration by parts, this is the Wallis product (Corollary~\ref{cor:wallis}) again.
\end{solution}

\begin{solution}{ex:quickexp}
\idea Put in the expansion of $H_n$ and keep every term down to $\frac1n$. Because $H_n$ is multiplied by $n+1$, a term of $H_n$ of order $\frac1{n^2}$ still contributes at order $\frac1n$.
\writeup Corollary~\ref{cor:harmonic} with $M=2$ gives $H_n=\ln n+\gamma+\frac1{2n}-\frac1{12n^2}+O(n^{-4})$. Multiply by $2(n+1)=2n+2$:
\[
\begin{aligned}2(n+1)H_n&=2n\ln n+2\gamma n+1-\frac1{6n}\eqbr+2\ln n+2\gamma+\frac1n+O\Bigl(\frac1{n^2}\Bigr).\end{aligned}
\]
Subtract $4n$ and collect terms. Then $-\frac1{6n}+\frac1n=\frac5{6n}$, which is the formula. The next term comes from $2\cdot\bigl(-\frac1{12n^2}\bigr)$: it is $-\frac1{6n^2}$.
\end{solution}

\begin{solution}{ex:zeta3}
\idea The proposition replaces the slow tail $\sum_{n>10}n^{-3}$, which is about $0.005$, by a few explicit terms and a remainder. The remainder is tiny because $\bar B_7$ is small and $x^{-10}$ falls fast.
\writeup With $s=3$, $(3)_{2k-1}=3\cdot4\cdots(2k+1)=\frac{(2k+1)!}2$, so the term $k$ is $\frac{B_{2k}}{(2k)!}\cdot\frac{(2k+1)!}2N^{-2-2k}=\frac{(2k+1)B_{2k}}2N^{-2k-2}$. For $k=1,2,3$ this is $\frac1{4N^4}$, $-\frac1{12N^6}$, $\frac1{12N^8}$. The other terms are $\frac{N^{-2}}2$ and $-\frac{N^{-3}}2$. With $N=10$ this is the formula, and $r=R_6$. Since $B_7=0$, the step from order $6$ to order $7$ in Theorem~\ref{thm:eulermaclaurin} adds nothing, so $R_6=R_7=\frac1{7!}\int_{10}^\infty\bar B_7(x)f^{(7)}(x)\dd x$ with $f^{(7)}(x)=-\frac{9!}2x^{-10}$. Lemma~\ref{lem:bernoullifn}(c) gives $|\bar B_7|\le\frac{2\cdot7!\,\zeta(7)}{(2\pi)^7}$, so
\[
\begin{aligned}|r|&\le\frac{2\zeta(7)}{(2\pi)^7}\cdot\frac{9!}2\int_{10}^\infty\frac{\mathrm dx}{x^{10}}\eqbr=\frac{2\zeta(7)}{(2\pi)^7}\cdot\frac{8!}{2\cdot10^9}\approx1.05\cdot10^{-10},\end{aligned}
\]
The bound uses $\zeta(7)\le1.01$. Numerically the first sum is $1.19753198567\ldots$. The formula and the true value are
\[
\begin{gathered}\zeta(3)\approx1.2020569031742,\qq\zeta(3)=1.2020569031596\ldots,\end{gathered}
\]
The actual error is $1.5\cdot10^{-11}$, seven times smaller than the bound.
\end{solution}

\begin{solution}{ex:partitionbound}
\idea The series $P(x)$ has positive terms, so any one of them, $p(n)x^n$, is at most $P(x)$. Bounding $P$ and choosing $x$ well is the saddle-point method used as an inequality.
\writeup (a) With $x=e^{-t}$, $\ln P(x)=\sum_k-\ln(1-x^k)=\sum_k\sum_{m\ge1}\frac{x^{km}}m$. The double series has positive terms, so the sum over $k$ may come first: $\sum_kx^{km}=\frac1{e^{mt}-1}\le\frac1{mt}$, since $e^u-1\ge u$. So $\ln P(e^{-t})\le\frac1t\sum_m\frac1{m^2}=\frac{\pi^2}{6t}$.

(b) For every $t>0$, $p(n)e^{-nt}\le P(e^{-t})$, so $\ln p(n)\le nt+\frac{\pi^2}{6t}$. The right side is smallest at $t=\frac\pi{\sqrt{6n}}$, and there it equals $2\pi\sqrt{n/6}=\pi\sqrt{2n/3}$.

(c) Let $m=\lfloor\sqrt n\rfloor$. For each subset $S$ of $\{2,\dots,m\}$, the sum of $S$ is at most $\frac{m(m+1)}2\le n$, because $m^2\le n$ and $m\le n$. Complete $S$ by $n-\sum S$ parts equal to $1$. That is a partition of $n$. Different subsets give different partitions: read off the parts larger than $1$. There are $2^{m-1}$ subsets.

The true size sits between the two bounds, and much closer to (b). Hardy and Ramanujan proved in 1918 that $p(n)\sim\frac{e^{\pi\sqrt{2n/3}}}{4\sqrt3\,n}$, so $\ln p(n)\sim\pi\sqrt{2n/3}$. Their proof studies $P$ near every root of unity, not only near $1$. That is the circle method. Chapter~\ref{ch:hr} proves the asymptotic formula. The ratio $\frac{\ln p(n)}{\pi\sqrt{2n/3}}$ creeps up: $0.74$ at $n=100$, $0.89$ at $n=1000$, $0.92$ at $n=2000$. The slow convergence comes from the factor $\frac1n$ in front.
\end{solution}

\begin{solution}{ex:motzkinasym}
\idea The singularity closest to $0$ is at $z=\frac13$, where $\sqrt{1-3z}$ vanishes. The other factor $\sqrt{1+z}$ is analytic there, and it only contributes its value.
\writeup For $n\ge1$, the polynomial $1-z$ does not contribute to $[z^{n+2}]$, so
\[
M_n=-\frac12[z^{n+2}]\sqrt{(1-3z)(1+z)}.
\]
Set $w=3z$. Then $[z^N]G(z)=3^N[w^N]G(w/3)$, and $\sqrt{(1-3z)(1+z)}=(1-w)^{1/2}h(w)$ with $h(w)=\sqrt{1+w/3}$, a power series of radius $3>1$. Theorem~\ref{thm:transfer}
with $\alpha=-\frac12$ and $\Gamma(-\frac12)=-2\sqrt\pi$ gives
\[
\begin{aligned}{}[w^N](1-w)^{1/2}h(w)&\sim\frac{h(1)N^{-3/2}}{\Gamma(-\frac12)}\eqbr=-\frac{1}{\sqrt3\,\sqrt\pi\,N^{3/2}},\end{aligned}
\]
because $h(1)=\frac2{\sqrt3}$. With $N=n+2\sim n$, one has $M_n\sim\frac12\cdot\frac{3^{n+2}}{\sqrt{3\pi}\,n^{3/2}}=\frac{3^{n+3/2}}{2\sqrt\pi\,n^{3/2}}$. At $n=200$ the ratio of the two sides is $0.988$.
\end{solution}

\begin{solution}{pb:unlabelled}
\idea Burnside writes $g_n$ as an average over $\mathfrak S_n$ of $2^{c(\sigma)}$, where $c(\sigma)$ is the number of cycles of $\sigma$ on pairs of vertices. The identity alone contributes $\frac{2^N}{n!}$. Every other permutation has far fewer cycles on pairs, and there are not enough of them to matter.

\writeup Let $\mathfrak S_n$ act on the $N$ pairs, and let $c(\sigma)$ be the number of cycles of $\sigma$ on pairs. A graph is a $2$-colouring of the pairs, so Burnside gives
\[
g_n=\frac1{n!}\sum_\sigma2^{c(\sigma)},\qquad\frac{n!\,g_n}{2^N}=1+\sum_{\sigma\ne e}2^{c(\sigma)-N}.
\]
Part (a) is immediate.

\emph{A bound on $c(\sigma)$.} Let $\sigma$ move exactly $m\ge2$ points. A pair fixed by $\sigma$ lies inside the $n-m$ fixed points, or is a $2$-cycle $\{a,\sigma(a)\}$ of $\sigma$. So there are at most $f=\binom{n-m}2+\frac m2$ fixed pairs. The other $N-f$ pairs lie in cycles of length at least $2$, so $c(\sigma)\le f+\frac{N-f}2$, and
\[
N-c(\sigma)\ge\frac{N-f}2=\frac12\Bigl(\binom n2-\binom{n-m}2-\frac m2\Bigr)=\frac{m(2n-m-2)}4.
\]

(b) and (c). There are at most $\binom nm m!\le n^m$ permutations moving $m$ points. A transposition $(a\,b)$ fixes the pair $\{a,b\}$ and the $\binom{n-2}2$ pairs that avoid $a$ and $b$, and it exchanges $\{a,x\}$ with $\{b,x\}$ for the $n-2$ other points $x$. So $N-c=n-2$ exactly, and the $\binom n2$ transpositions contribute $\binom n22^{2-n}=\frac{n(n-1)}{2^{n-1}}$. For $m=3$, those $\sigma$ contribute at most $n^32^{-3(2n-5)/4}=O(n^32^{-3n/2})$. For $4\le m\le\frac n2$, one has $2n-m-2\ge\frac{3n-4}2$, so they contribute at most $\bigl(n2^{-(3n-4)/8}\bigr)^m$. That is a geometric series in $m$, and the sum over $m\ge4$ is $O(n^42^{-3n/2})$. For $m>\frac n2$, one has $2n-m-2\ge n-2$, so $N-c\ge\frac{n(n-2)}8$, and the at most $n!$ such $\sigma$ contribute at most $n^n2^{-n(n-2)/8}$, which is far smaller. Altogether
\[
\frac{n!\,g_n}{2^N}=1+\frac{n(n-1)}{2^{n-1}}+O\bigl(n^42^{-3n/2}\bigr)\longrightarrow1.
\]
\tools Volume~3, Chapter~3 (Burnside's lemma). The rest is geometric series.
\leads The statement $g_n\sim\frac{2^N}{n!}$ means that almost every graph has a trivial automorphism group (Erdős and Rényi, 1963). Oberschelp (1967) found the full asymptotic expansion. The ratio converges quickly: at $n=10$ it is already $1.24$, and the two-term formula predicts $1.18$.
\end{solution}

\begin{solution}{pb:oddcycles}
\idea Restrict the cycle lengths to odd numbers or to even numbers. That gives $\exp\sum_{j\text{ odd}}\frac{x^j}j=\sqrt{\frac{1+x}{1-x}}$ and $\exp\sum_{j\text{ even}}\frac{x^j}j=\frac1{\sqrt{1-x^2}}$, and $\sqrt{\frac{1+x}{1-x}}=\frac{1+x}{\sqrt{1-x^2}}$. The even parts agree.

\writeup (a) The exponential formula, as in Problem~\ref{pb:roots}, gives
\[
\sum o_n\frac{x^n}{n!}=\exp\Bigl(\sum_{j\text{ odd}}\frac{x^j}j\Bigr)=\exp\Bigl(\frac12\log\frac{1+x}{1-x}\Bigr)=\frac{1+x}{\sqrt{1-x^2}},
\]
and
\[
\sum e_n\frac{x^n}{n!}=\exp\Bigl(\sum_{j\text{ even}}\frac{x^j}j\Bigr)=\exp\Bigl(-\frac12\log(1-x^2)\Bigr)=\frac1{\sqrt{1-x^2}}.
\]
Now $\frac1{\sqrt{1-x^2}}=\sum_m\binom{2m}m\frac{x^{2m}}{4^m}$ (Proposition~\ref{prop:binomcoef} with $\alpha=\frac12$, where $a_m(\frac12)=\binom{2m}m4^{-m}$). So for $n=2m$,
\[
o_n=e_n=(2m)!\binom{2m}m4^{-m}=\Bigl(\frac{(2m)!}{2^mm!}\Bigr)^2=\bigl(1\cdot3\cdots(2m-1)\bigr)^2,
\]
The coefficient of $x^{2m+1}$ in $\frac{1+x}{\sqrt{1-x^2}}$ is $\binom{2m}m4^{-m}$ as well, so $o_{2m+1}=(2m+1)!\binom{2m}m4^{-m}=(2m+1)o_{2m}$.

(b) The order of $\sigma$ is the least common multiple of its cycle lengths. That order is odd exactly when every cycle is odd. The probability is $\frac{o_n}{n!}=\binom{2m}m4^{-m}$ with $m=\lfloor\frac n2\rfloor$, and $\binom{2m}m4^{-m}\sim\frac1{\sqrt{\pi m}}$ (Example~\ref{ex:centralsaddle}), which is $\sim\sqrt{\frac2{\pi n}}$.
\tools Chapter~\ref{ch:generating} (the exponential formula); Proposition~\ref{prop:binomcoef} and Example~\ref{ex:centralsaddle} of this chapter.
\leads For even $n$ the equality $o_n=e_n$ asks for a bijection between permutations with all cycles odd and permutations with all cycles even. Several are known. None of them is obvious. Finding one is a good challenge.
\end{solution}

\begin{solution}{pb:powersum}
\idea The terms $(1-\frac jn)^n$ are close to $e^{-j}$. The sum is dominated by the last few values of $k$, and the correction comes from the second-order term of $\ln(1-x)$.

\writeup For $0\le x\le\frac12$, $\ln(1-x)=-\sum_{i\ge1}\frac{x^i}i$, so $\ln(1-x)\le-x-\frac{x^2}2$ and $\ln(1-x)\ge-x-\frac{x^2}2-\frac{x^3}3\sum_{i\ge0}x^i\ge-x-\frac{x^2}2-x^3$. Set $a_j=(1-\frac jn)^n$, so the sum is $\sum_{j=0}^{n-1}a_j$. For $j\le\frac n2$, with $x=\frac jn$,
\[
e^{-j}e^{-j^2/(2n)}e^{-j^3/n^2}\le a_j\le e^{-j}e^{-j^2/(2n)}.
\]
Since $1-y\le e^{-y}\le1-y+\frac{y^2}2$ for $y\ge0$, both sides equal $e^{-j}\bigl(1-\frac{j^2}{2n}+E_j\bigr)$ with $|E_j|\le\frac{j^4}{8n^2}+\frac{j^3}{n^2}$. So
\[
\sum_{j\le n/2}a_j=\sum_{j\le n/2}e^{-j}-\frac1{2n}\sum_{j\le n/2}j^2e^{-j}+O\Bigl(\frac1{n^2}\sum_jj^4e^{-j}\Bigr).
\]
Replacing the first two sums by the full series costs $O(n^2e^{-n/2})$. For $\frac n2<j<n$, one has $a_j\le2^{-n}$, so those terms contribute at most $n2^{-n}$. With $\sum_je^{-j}=\frac e{e-1}$ and $\sum_jj^2x^j=\frac{x(1+x)}{(1-x)^3}$ at $x=\frac1e$, which equals $\frac{e(e+1)}{(e-1)^3}$, the claimed expansion follows.

\tools{Volume~1, Chapters~3 and~8 (series, the expansion of $\ln(1-x)$); this chapter (the cutting used in Laplace's method).}
\leads Compare with $\sum_kk^k\sim n^n$ (Exercise~\ref{ex:kk}). There only the last term matters. Here the last $O(1)$ terms all count. The method cuts the sum into a region where a Taylor expansion is uniform and a region where everything is exponentially small. That is the discrete Laplace method, the same cut that proves Theorem~\ref{thm:laplace}.
\end{solution}

\begin{solution}{pb:birthdaytriple}
\idea This is Laplace's method with a twist. The integrand is the $n$-th power of a function whose maximum sits at the endpoint $0$, where it is flat to third order. The window that matters has width $n^{-1/3}$ instead of $n^{-1/2}$, and the Gaussian integral is replaced by $\int_0^\infty e^{-y^3}\dd y$.

\writeup (a) Put $t=nx$. Then $\mathbb E[T]=n\int_0^\infty\psi(x)^n\dd x$ with $\psi(x)=e^{-x}(1+x+\frac{x^2}2)$. Now $\psi(0)=1$ and $\psi'(x)=-\frac{x^2}2e^{-x}<0$ for $x>0$, so $\psi$ decreases. Also $\psi(x)=1-e^{-x}\sum_{k\ge3}\frac{x^k}{k!}=1-\frac{x^3}6+O(x^4)$, so $\ln\psi(x)=-\frac{x^3}6+O(x^4)$ near $0$.

Fix $\varepsilon\in\left]0,1\right[$ and choose $\delta>0$ so that
\[
-(1+\varepsilon)\frac{x^3}6\le\ln\psi(x)\le-(1-\varepsilon)\frac{x^3}6\qquad(0\le x\le\delta).
\]
\emph{The tail.} On $[\delta,\infty[$, one has $\psi^n\le\psi(\delta)^{n-1}\psi$, and $\int_0^\infty\psi=1+1+1=3$. So $n\int_\delta^\infty\psi^n\le3n\,\psi(\delta)^{n-1}$, which tends to $0$.

\emph{The window.} For $c>0$, the substitution $x=(\frac6{cn})^{1/3}y$ produces
\[
n\int_0^\delta e^{-cnx^3/6}\dd x=\Bigl(\frac{6n^2}c\Bigr)^{1/3}\int_0^{Y_n}e^{-y^3}\dd y,\qquad Y_n=\delta\Bigl(\frac{cn}6\Bigr)^{1/3}\to\infty.
\]
And $\int_0^\infty e^{-y^3}\dd y=\frac13\int_0^\infty u^{-2/3}e^{-u}\dd u=\frac13\Gamma(\frac13)=\Gamma(\frac43)$ (substitute $u=y^3$; Chapter~\ref{ch:gamma}). Take $c=1-\varepsilon$ for an upper bound and $c=1+\varepsilon$ for a lower bound. Then
\[
(1+\varepsilon)^{-1/3}\le\liminf\frac{\mathbb E[T]}{\Gamma(\frac43)(6n^2)^{1/3}}\le\limsup\frac{\mathbb E[T]}{\Gamma(\frac43)(6n^2)^{1/3}}\le(1-\varepsilon)^{-1/3}.
\]
Let $\varepsilon\to0$.

(b) The approximation is $\Gamma(\frac43)(6\cdot365^2)^{1/3}=82.9$, against the exact $88.74$. That is seven per cent too low. For pairs, the same method gives $\sqrt{\pi\cdot365/2}=23.9$ against $24.62$.

\tools{Chapter~\ref{ch:generating} (the integral for $\mathbb E[T]$); Chapter~\ref{ch:gamma} (the Gamma function); this chapter (the cutting argument of Laplace's method).}
\leads Klamkin and Newman (1967) found the general answer for a value seen $m$ times: $\mathbb E[T]\sim(m!)^{1/m}\Gamma(1+\frac1m)\,n^{1-1/m}$. For $m=2$ this is $\sqrt{\pi n/2}$. The relative error decays only like $n^{-1/3}$ for $m=3$, which is why $82.9$ is still some way from $88.74$.
\end{solution}

\begin{solution}{pb:sqrtsum}
\idea Proposition~\ref{prop:zetaem} is an identity between holomorphic functions, valid far to the left of $\operatorname{Re}s=1$. Read at $s=-\frac12$, it is Euler--Maclaurin for $\sum\sqrt k$, with $\zeta(-\frac12)$ playing the role of the constant.

\writeup Proposition~\ref{prop:zetaem} says that for $N,M\ge1$ and $\operatorname{Re}s>1-2M$, $s\ne1$,
\[
\zeta(s)=\sum_{k\le N}k^{-s}+\frac{N^{1-s}}{s-1}-\frac{N^{-s}}2+\sum_{j=1}^M\frac{B_{2j}}{(2j)!}(s)_{2j-1}N^{1-s-2j}-\frac{(s)_{2M}}{(2M)!}\int_N^\infty\frac{\bar B_{2M}(x)}{x^{s+2M}}\dd x,
\]
Here $(s)_i=s(s+1)\cdots(s+i-1)$ and $\bar B_{2M}$ is bounded. Take $M=2$. For $s=-\frac12$: $\frac{N^{3/2}}{-3/2}=-\frac23N^{3/2}$, $-\frac{N^{1/2}}2$, and the $j=1$ term is $\frac1{12}\cdot(-\frac12)N^{-1/2}=-\frac1{24\sqrt N}$. The $j=2$ term and the integral are $O(N^{-5/2})$. So $\sum_{k\le N}\sqrt k=\zeta(-\frac12)+\frac23N^{3/2}+\frac12\sqrt N+\frac1{24\sqrt N}+O(N^{-5/2})$. For $s=\frac12$: $\frac{N^{1/2}}{-1/2}=-2\sqrt N$, $-\frac1{2\sqrt N}$, and the $j=1$ term is $\frac1{12}\cdot\frac12N^{-3/2}=\frac1{24N^{3/2}}$. The rest is $O(N^{-7/2})$. So $\sum_{k\le N}\frac1{\sqrt k}=\zeta(\frac12)+2\sqrt N+\frac1{2\sqrt N}-\frac1{24N^{3/2}}+O(N^{-7/2})$.

\tools{Theorem~\ref{thm:eulermaclaurin} (Euler--Maclaurin) and Proposition~\ref{prop:zetaem} (the continuation of $\zeta$).}
\leads One has $\zeta(\frac12)<0$ even though every term of the series is positive. The constant is what remains after $2\sqrt N$ is removed. It is not itself a sum. The functional equation (Corollary~\ref{cor:zetaplane}(d)) gives $\zeta(-\frac12)=-\frac{\zeta(3/2)}{4\pi}$. Ramanujan called such constants ``the constant of a series''. Hardy's \emph{Divergent Series} (1949) makes the idea systematic.
\end{solution}

\begin{solution}{ex:rw11a}
\idea Stop after the term $k=3$. The next remainder is positive and strictly smaller than the first omitted term of the majorant series.
\writeup The partial sum through $k=3$ is $1-\frac15+\frac{2}{25}-\frac{6}{125}=\frac{125-25+10-6}{125}=\frac{104}{125}=0.832$. Proposition~\ref{prop:eulerseries} with $K=4$ says the remainder equals $+(1/5)^4\int_0^\infty t^4e^{-t}/(1+t/5)\dd t$, so it is positive, and the integral is strictly smaller than $4!=24$ because $1+t/5>1$ on a set of positive length. Hence $0<r_4<24/625$ and $\frac{104}{125}<f\bigl(\frac15\bigr)<\frac{104}{125}+\frac{24}{625}=\frac{544}{625}=0.8704$. The integral itself equals $0.8521108814\ldots$, which sits inside that interval.
\end{solution}

\begin{solution}{ex:rw11b}
\idea Write the four explicit terms at $n=5$ and compare with the exact fraction $137/60$. The corollary names the sign of what you left out.
\writeup The expansion through $\frac1{120n^4}$ is $\ln5+\gamma+\frac1{10}-\frac1{12\cdot25}+\frac1{120\cdot625}=\ln5+\gamma+\frac1{10}-\frac1{300}+\frac1{75000}=2.283333577\ldots$. And $H_5=\frac{137}{60}=2.283333333\ldots$. The approximation is larger by $2.440\cdot10^{-7}$. The first omitted term is $-\frac1{252\cdot5^6}=-\frac1{3937500}=-2.540\cdot10^{-7}$. The error has that sign, and $2.440\cdot10^{-7}<2.540\cdot10^{-7}$.
\end{solution}

\begin{solution}{ex:rw11c}
\idea An endpoint maximum with a non-zero derivative produces a spike of width $1/n$. An interior maximum with $\varphi''<0$ produces a bump of width $1/\sqrt n$, as in Figure~\ref{fig:rw11-bump}.
\writeup The antiderivative of $e^{-nx}$ is $-e^{-nx}/n$, so $\int_0^\infty e^{-nx}\dd x=\frac1n$. For the Gaussian, set $u=x\sqrt n$. Then $\int_{-\infty}^\infty e^{-nx^2}\dd x=n^{-1/2}\int_{-\infty}^\infty e^{-u^2}\dd u=\sqrt{\pi/n}$. At $n=10$ the endpoint integral equals $0.1$, and the Gaussian integral equals $\sqrt{\pi/10}=0.560499122\ldots$. The interior bump is the larger of the two.
\end{solution}

\begin{solution}{ex:rw11d}
\idea Compute the leading saddle approximation first, then multiply by the factor $1-1/(8n)$ from the transfer of $(1-4z)^{-1/2}$.
\writeup For $n=2$, $4^2/\sqrt{2\pi}=16/\sqrt{2\pi}=6.383076486\ldots$, against $\binom42=6$. The leading term is high by about $6.4\%$. For $n=10$, $4^{10}/\sqrt{10\pi}=187078.9729\ldots$, against $\binom{20}{10}=184756$. The two-term approximation multiplies by $1-\frac1{80}=\frac{79}{80}$ and equals $184740.4858\ldots$. It misses $184756$ by $15.514$, a relative error of $8.4\cdot10^{-5}$.
\end{solution}

\section*{Sources and further reading}

\begin{itemize}
\item N.~H.~Abel, letter to B.~M.~Holmboe, 16 January 1826, in \emph{\OE uvres compl\`etes}, ed.\ L.~Sylow and S.~Lie, Christiania, 1881,
vol.~2, 256--258. The epigraph is the usual English rendering.
\item L.~Euler, \emph{Inventio summae cuiusque seriei ex dato termino generali} (E47), Commentarii academiae scientiarum Petropolitanae 8
(1736, published 1741), 9--22: the summation formula and the twenty decimals of $\sum\frac1{n^2}$. \emph{De seriebus divergentibus} (E247),
Novi commentarii 5 (1754/55, published 1760), 205--237: the value $0.596$ of $0!-1!+2!-\cdots$ C.~Maclaurin, \emph{A Treatise of Fluxions}, Edinburgh, 1742.
\item P.~Erd\H os and A.~R\'enyi, \emph{Asymmetric graphs}, Acta Mathematica Academiae Scientiarum Hungaricae 14 (1963), 295--315.
M.~S.~Klamkin and D.~J.~Newman, \emph{Extensions of the birthday surprise}, Journal of Combinatorial Theory 3 (1967), 279--282.
\item N.~G.~de~Bruijn, \emph{Asymptotic Methods in Analysis}, North-Holland, 1958 (Dover reprint 1981). The classic: Laplace's method, saddle points,
Euler--Maclaurin, iterated functions, all through worked examples.
\item R.~L.~Graham, D.~E.~Knuth and O.~Patashnik, \emph{Concrete Mathematics}, 2nd ed., Addison-Wesley, 1994, Chapter~9 (asymptotics, Euler's
summation formula, Stirling's series).
\item P.~Flajolet and R.~Sedgewick, \emph{Analytic Combinatorics}, Cambridge University Press, 2009, Chapters~VI and~VIII (singularity analysis and the
saddle point method). Free online from the authors.
\item P.~Flajolet and A.~Odlyzko, \emph{Singularity analysis of generating functions}, SIAM Journal on Discrete Mathematics 3 (1990), 216--240.
\item E.~T.~Whittaker and G.~N.~Watson, \emph{A Course of Modern Analysis}, Cambridge, 1927, Chapter~12 (Stirling's series with the bound of
Theorem~\ref{thm:stirlingseries}).
\item G.~H.~Hardy, \emph{Divergent Series}, Oxford, 1949. What divergent series can mean, by the analyst who made them respectable; Chapter~1 opens
with Abel's letter.
\item D.~J.~Newman, \emph{Analytic Number Theory}, Springer, 1998, Chapter~I (a short proof of the Hardy--Ramanujan asymptotic for $p(n)$).
\end{itemize}
